
% ============================================================
% HANDOFF FOR HIMANSHU
% Resolve the mismatch between the theoretical rank mapping
% and the three implemented ranking strategies.
%
% INSERTION INSTRUCTIONS
%
% 1. Replace the current Section II.D, "Problem Formulation,"
%    in full with the section below.
%
% 2. In Section IV.B, "Workflow Primitives," replace the current
%    Table I with the revised table below.
%
% 3. Immediately after the revised Table I, insert the short
%    validation paragraph below.
%
% 4. Use the following implementation rule throughout the paper:
%    retain each strategy-specific prediction for evaluation,
%    compute the deterministic theory-implied rank separately,
%    and flag any disagreement. Do not silently overwrite the
%    strategy-specific prediction.
% ============================================================


% ============================================================
% SUBMISSION PRIORITY-- JULY 21, 2026 DUE DATE
%
% First: replace Section II.D and Table I; insert the revised figure.
% Second: reconcile sample sizes and populate or remove all TBD tables.
% Third: apply the Priority A language changes and anonymize the paper.
% Apply Priority B terminology changes only after those tasks are complete.
% ============================================================


% ============================================================
% ADD THIS POSITIONING NOTE NEAR THE TOP OF THE HANDOFF
% ============================================================
%
% POSITIONING HIERARCHY
%
% Keep "theory-guided" in the paper and title, but make
% replicability the primary contribution.
%
% Primary contribution:
% A replicable LLM-based measurement workflow.
%
% Distinctive design principle:
% The established theory of delegation and discretion determines
% the workflow variables, stage ordering, and consistency rules.
%
% Theory-guided should not be presented as a separate new theory
% contribution. It explains why the workflow is substantively
% structured rather than an arbitrary chain of prompts.
% ============================================================


% ============================================================
% RECOMMENDED TITLE
% ============================================================
%
% A Replicable Theory-Guided LLM Workflow for Measuring Legal
% Delegation and Administrative Discretion
% ============================================================


% ============================================================
% ADD OR REPLACE THE CENTRAL CONTRIBUTION SENTENCE IN THE
% INTRODUCTION
% ============================================================

This paper develops a replicable LLM-based measurement instrument
in which an established theory of delegation and discretion
determines the workflow's variables, ordering, and consistency rules.


% ============================================================
% ADD THIS SHORT CLARIFICATION AT THE END OF THE INTRODUCTION,
% IMMEDIATELY BEFORE THE CONTRIBUTION LIST
% ============================================================

The principal contribution is replicability. The theory-guided
design makes the workflow substantively meaningful by ensuring that 
its stages and dependencies align with the established measurement
model rather than with an arbitrary sequence of prompts.


% ============================================================
% REVISE THE CONTRIBUTION LANGUAGE SO THAT REPLICABILITY COMES
% FIRST AND THEORY-GUIDED DESIGN EXPLAINS THE ARCHITECTURE
% ============================================================

\begin{enumerate}
    \item The paper develops a replicable LLM-based measurement
    instrument for legal delegation and administrative discretion.
    The complete workflow, rather than an isolated prompt, is
    archived, versioned, and auditable.

    \item The workflow is theory-guided: the established measurement
    model determines the variables, stage ordering, conditional
    dependencies, and permitted relationships between the component
    classifications and the final rank.

    \item The system separates generative LLM judgments from
    deterministic theoretical mappings and preserves the intermediate
    evidence needed to diagnose disagreements.

    \item The prototype enables domain researchers to apply and inspect
    the measurement process without directly managing prompts, APIs,
    model orchestration, or structured-output parsing.
\end{enumerate}


% ============================================================
% ADD THIS LANGUAGE TO THE DISCUSSION SECTION
% ============================================================

Our principal contribution is replicability. The theory-guided design
makes the replicable workflow substantively meaningful by ensuring
that its stages and dependencies align with the established
measurement of a discretion model \citep{epstein1999delegating}. 
The theoretical framework provides the formal structure that the workflow implements and preserves.


% ============================================================
% ADD THIS LANGUAGE TO THE NOVELTY STATEMENT
% ============================================================

The paper relies on well-established theories of delegation or
administrative discretion. Our contribution is a replicable LLM
measurement architecture that implements a theoretical model through
explicit variables, ordered stages, deterministic consistency rules,
and preserved execution records.


% ============================================================
% ADD THIS LANGUAGE TO THE CONCLUSION
% ============================================================

The resulting system is best understood as a replicable measurement
instrument, with its computational structure guided by an established
theory of delegation and discretion. Replicability is the primary
contribution; theory-guided design explains why the workflow's stages,
dependencies, and validation rules are substantively meaningful.


% ============================================================
% REPLACEMENT FOR SECTION II.D
% ============================================================

\subsection{Problem Formulation}

Let \(x_i\) denote the legal text for law \(i\). The theoretical
measurement model contains three component variables. Let

\[
d_i \in \{0,1\}
\]

indicate whether the law delegates qualifying policymaking authority,

\[
a_i \in \{\mathrm{Low},\mathrm{High}\}
\]

denote the amount of delegated authority, and

\[
c_i \in \{\mathrm{Strong},\mathrm{Weak}\}
\]

denote the strength of statutory constraints.

The substantive theory defines a deterministic mapping from these
components to the discretion rank:

\[
y_i = m(d_i,a_i,c_i),
\]

where

\[
m(d_i,a_i,c_i)=
\begin{cases}
0, & d_i=0,\\
1, & d_i=1,\ a_i=\mathrm{Low},\ c_i=\mathrm{Strong},\\
2, & d_i=1,\ a_i=\mathrm{Low},\ c_i=\mathrm{Weak},\\
3, & d_i=1,\ a_i=\mathrm{High},\ c_i=\mathrm{Strong},\\
4, & d_i=1,\ a_i=\mathrm{High},\ c_i=\mathrm{Weak}.
\end{cases}
\]

This mapping defines the substantive meaning of the five discretion
ranks. It is part of the measurement model rather than a learned
relationship.

The workflow uses versioned LLM stages to estimate the component
variables:

\[
\hat d_i = g_D(x_i),
\]

\[
\hat a_i = g_A(x_i,\hat d_i),
\]

\[
\hat c_i = g_C(x_i,\hat d_i,\hat a_i).
\]

The component stages also return structured evidence and rationales,
denoted collectively by \(r_i\). Their prompts, inputs, outputs,
evidence, and model settings are retained in the execution trace.

The system then applies one of three rank-classification strategies.
Let

\[
\mathcal{S}
=
\{\mathrm{Cascade},\mathrm{Multiclass},\mathrm{TwoBand}\}
\]

denote the Sequential Ordinal Cascade, Direct Multiclass Rank
Prediction, and Two-Band Rank Classification strategies, respectively. For strategy
\(s \in \mathcal{S}\),

\[
\tilde y_i^{(s)}
=
h_s(x_i,\hat d_i,\hat a_i,\hat c_i,r_i).
\]

The Sequential Ordinal Cascade exploits the outcome's ordered structure. The delegation gate first separates Rank \(0\) from Ranks \(1\)--\(4\). Conditional on delegation, the cascade proceeds in stages, with the first separating Rank \(1\) from Ranks \(2\)--\(4\), the second separating Rank \(2\) from Ranks \(3\)--\(4\), and the final stage separating Rank \(3\) from Rank \(4\). Each positive decision produces an early exit, while each negative decision routes the case to the next stage.

The Direct Multiclass strategy predicts one of Ranks \(1\)--\(4\) in
a single ranking step using the stored component outputs and
rationales. The Two-Band strategy first separates the lower-rank band
\(\{1,2\}\) from the higher-rank band \(\{3,4\}\), then classifies
within the selected band.

The theoretical mapping and the ranking strategies, therefore, perform
different functions. The mapping \(m(d_i,a_i,c_i)\) defines the
construct, whereas the functions \(h_s(\cdot)\) estimate the final
rank under alternative computational strategies.

For each law, the system also computes the theory-implied rank; that is, the rank is determined by the estimated component classifications,

\[
y_i^{\mathrm{map}}
=
m(\hat d_i,\hat a_i,\hat c_i)
\]

and compares it with the strategy-specific prediction. Define the
consistency indicator

\[
q_i^{(s)}
=
\mathbf{1}
\left[
\tilde y_i^{(s)}
=
y_i^{\mathrm{map}}
\right].
\]

A value of \(q_i^{(s)}=1\) indicates that the rank predicted by
strategy \(s\) agrees with the rank implied by the model's own
delegation, authority, and constraint classifications. A disagreement
is retained and flagged for review rather than silently overwritten.
This design allows the evaluation to distinguish benchmark agreement
from internal theoretical consistency.


% ============================================================
% REPLACEMENT FOR TABLE I IN SECTION IV.B
% ============================================================

\begin{table}[t]
\caption{Algorithmic execution of the theory-guided discretion workflow.}
\label{tab:algorithm}
\centering
\footnotesize
\begin{tabular}{p{0.06\columnwidth}p{0.88\columnwidth}}
\hline
\textbf{Step} & \textbf{Operation and Logic} \\
\hline

1 &
\textbf{Load Source Text:}
Read the legislative summary or statutory text \(x_i\) and attach
stable identifiers. \\

2 &
\textbf{Delegation Gate (LLM):}
Estimate
\(\hat d_i=g_D(x_i)\in\{0,1\}\)
and retain the supporting evidence and rationale. \\

3 &
\textbf{Early Exit (Deterministic):}
If \(\hat d_i=0\), set the theory-implied rank
\(y_i^{\mathrm{map}}=0\), assign the required default values, and stop
the delegated-discretion branch. \\

4 &
\textbf{Component Classification (LLM):}
Conditional on \(\hat d_i=1\), identify the relevant administrative
actors, classify delegated authority
\(\hat a_i\in\{\mathrm{Low},\mathrm{High}\}\), classify statutory
constraints
\(\hat c_i\in\{\mathrm{Strong},\mathrm{Weak}\}\), and retain the
supporting evidence and rationales \(r_i\). \\

5 &
\textbf{Strategy-Specific Rank Prediction (LLM):}
Apply the selected rank-classification strategy. The Sequential
Ordinal Cascade uses
successive early-exit binary decisions; Direct Multiclass predicts
Ranks \(1\)--\(4\) in one step; and Two-Band first predicts a
lower/higher rank band and then classifies within the selected band. \\

6 &
\textbf{Theoretical Consistency Validation (Deterministic):}
Compute
\(y_i^{\mathrm{map}}=m(\hat d_i,\hat a_i,\hat c_i)\)
and compare it with the strategy-specific prediction
\(\tilde y_i^{(s)}\). Record the consistency indicator
\(q_i^{(s)}\) and flag any disagreement. \\

7 &
\textbf{Output Serialization:}
Store the component classifications, theory-implied rank,
strategy-specific rank, consistency flag, evidence, prompts, model
settings, execution trace, runtimes, token use, cost, and final
reported output. \\

\hline
\end{tabular}
\end{table}


% ============================================================
% INSERT IMMEDIATELY AFTER TABLE I
% ============================================================

The workflow retains both the strategy-specific prediction
\(\tilde y_i^{(s)}\) and the theory-implied rank
\(y_i^{\mathrm{map}}\). Strategy-specific predictions are used for
comparative evaluation and are not silently replaced when they conflict
with the deterministic component mapping. Instead, the validation node
records the disagreement with an inconsistency flag. This separation allows 
researchers to evaluate both agreement with the expert-coded benchmark and 
consistency with the rules used to construct the discretion rank.


% ============================================================
% ONE GLOBAL TERMINOLOGY CHANGE
% ============================================================
%
% Replace statements saying that "the final rank is deterministic"
% with:
%
% "The theoretical rank mapping is deterministic, while the final
% strategy-specific rank is produced by one of three alternative
% rank-classification strategies and validated against the theory-implied rank."
%
% This distinction should be reflected in the Abstract, Introduction,
% Figure captions, and Conclusion where necessary.


% ============================================================
% FIGURE 1: Four-stage theory-guided classification workflow
%
% Required preamble packages/libraries:
%   \usepackage{tikz}
%   \usetikzlibrary{arrows.meta,positioning,calc}
%   \usepackage{tabularx}
%   \usepackage{array}
%
% IEEEtran should control caption formatting. Do not load the
% caption package solely for this figure.
% ============================================================

\begin{figure*}[t]
\centering

% ---- Four-step overview ----
\begin{tikzpicture}[
  font=\small\sffamily,
  stage/.style={
    draw,
    rounded corners,
    very thick,
    align=center,
    minimum height=1.05cm,
    text width=3.25cm,
    inner sep=5pt
  },
  stagearrow/.style={-{Latex[length=3mm,width=2mm]}, very thick},
  node distance=0.52cm
]

\node[stage] (s1) {\textbf{Step 1}\\Input};
\node[stage, right=of s1] (s2) {\textbf{Step 2}\\Delegation Gate};
\node[stage, right=of s2, text width=4.35cm] (s3)
{\textbf{Step 3}\\Component Classification and Rank Assignment};
\node[stage, right=of s3] (s4) {\textbf{Step 4}\\Validation and Storage};

\draw[stagearrow] (s1) -- (s2);
\draw[stagearrow] (s2) -- (s3);
\draw[stagearrow] (s3) -- (s4);

\end{tikzpicture}

\vspace{0.7em}

% ---- Detailed workflow ----
\begin{tikzpicture}[
  font=\small\sffamily,
  box/.style={
    draw,
    rounded corners,
    thick,
    align=center,
    minimum height=0.95cm,
    text width=2.45cm,
    inner sep=4pt
  },
  branchbox/.style={
    draw,
    rounded corners,
    thick,
    align=center,
    minimum height=1.0cm,
    text width=2.75cm,
    inner sep=4pt
  },
  arrow/.style={-{Latex[length=2.2mm]}, thick},
  node distance=0.72cm and 0.82cm
]

\node[box] (text) {Legal Text};
\node[box, right=of text] (gate) {Delegation Gate};
\node[box, below=0.85cm of gate] (exit) {Deterministic Exit\\Rank 0};
\node[box, right=of gate] (components) {Component Classification};

\node[branchbox, above right=0.62cm and 0.72cm of components] (mapping)
{Theory-Based Rank\\Mapping};
\node[branchbox, below right=0.62cm and 0.72cm of components] (strategy)
{Rank-Classification\\Strategy};

\node[box, right=1.0cm of components] (check) {Consistency Check};
\node[box, right=of check] (output) {Coded Output};

\draw[arrow] (text) -- (gate);
\draw[arrow] (gate) -- node[above, font=\footnotesize\sffamily] {Yes} (components);
\draw[arrow] (gate) -- node[right, font=\footnotesize\sffamily] {No} (exit);
\draw[arrow] (components) -- (mapping);
\draw[arrow] (components) -- (strategy);
\draw[arrow] (mapping) -- (check);
\draw[arrow] (strategy) -- (check);
\draw[arrow] (check) -- (output);

\end{tikzpicture}

\vspace{0.55em}

% ---- Variables and records by step ----
\begin{minipage}{0.97\textwidth}
\centering
\footnotesize
\textbf{Variables and records by step}

\vspace{0.35em}

\begin{tabularx}{\textwidth}{@{}>{\raggedright\arraybackslash}X@{\hspace{1.2em}}>{\raggedright\arraybackslash}X@{}}

\textbf{Step 1. Input:}
A legislative summary or statutory text, together with a stable law identifier.
&
\textbf{Step 2. Delegation Gate:}
$\hat d_i\in\{0,1\}$, indicating whether the law contains a qualifying delegation.
\\[0.55em]

\textbf{Step 3. Component Classification and Rank Assignment:}
$\hat a_i\in\{\mathrm{Low},\mathrm{High}\}$,
$\hat c_i\in\{\mathrm{Strong},\mathrm{Weak}\}$, and $r_i$, containing supporting evidence and rationale. The workflow records both
$y_i^{\mathrm{map}}=m(\hat d_i,\hat a_i,\hat c_i)$, the rank implied by the component classifications, and $\tilde y_i^{(s)}$, the rank predicted by strategy $s$.
&
\textbf{Step 4. Validation and Storage:}
$q_i^{(s)}=\mathbf{1}[\tilde y_i^{(s)}=y_i^{\mathrm{map}}]$, together with the final rank, component values, evidence, model settings, and execution trace.
\\

\end{tabularx}
\end{minipage}

\caption{Four-step theory-guided classification workflow. Step 1 ingests a legislative summary or statutory text and a stable law identifier. Step 2 determines whether the law contains a qualifying delegation and assigns Rank 0 when delegation is absent. Step 3 classifies authority and constraints and produces both the theory-based and strategy-predicted ranks. Step 4 compares the two rank outputs and stores the coded variables, evidence, settings, and execution trace.}
\label{fig:theory_guided_workflow}
\end{figure*}


% ============================================================
% OPTIONAL SENTENCE TO ADD IMMEDIATELY BEFORE THE FIGURE
% ============================================================
%
% Figure~\ref{fig:theory_strategy_separation} separates the
% deterministic theoretical mapping from the three implemented
% rank-classification strategies and shows how the workflow records internal
% consistency without silently overwriting strategy-specific
% predictions.
% ============================================================


% ============================================================
% PAPER-WIDE LANGUAGE AND TERMINOLOGY REVIEW
% Apply this as a focused consistency pass after inserting the
% formal-model replacement above. These are not optional stylistic
% preferences: they remove overclaims, undefined terms, and internal
% inconsistencies that technical reviewers are likely to notice.
% ============================================================

% ============================================================
% PRIORITY A: REQUIRED BEFORE SUBMISSION
% ============================================================

% ------------------------------------------------------------
% TITLE
% ------------------------------------------------------------
% KEEP:
% A Replicable Theory-Guided LLM Workflow for Measuring Legal
% Delegation and Administrative Discretion
%
% RATIONALE:
% Replicability is the primary contribution. Theory-guided describes
% the design principle that determines the variables, ordering, and
% consistency rules. Do not present theory-guided AI as a new theory.


% ------------------------------------------------------------
% ABSTRACT, PAGE 1
% ------------------------------------------------------------
% 1. Replace "large-language-model" with "large language model."
%
% 2. Replace the sentence that currently states that all three ranking
%    strategies are evaluated across six models against a benchmark of
%    N=169. The delegation and discretion tasks may have different
%    samples, and the final counts must match the completed runs.
%
% USE AFTER THE SAMPLE AUDIT AND RUNS ARE COMPLETE:

We evaluate the Delegation Gate on the full delegation benchmark and
compare three rank-classification strategies across multiple language
models on the subset with complete expert-coded discretion ranks.

% Add the exact sample sizes and headline results only after Tables IV--
% VIII have been populated and checked.
%
% 3. Replace "ensures logical consistency and strict reproducibility."
%
% USE:

By separating generative judgments from deterministic theoretical
mappings and preserving the execution record, the system enforces
specified consistency rules and supports reproducible auditing and
rerunning.

% 4. Replace or revise the Index Terms so that they identify the NLP
%    contribution directly.
%
% SUGGESTED INDEX TERMS:
% legal natural language processing, large language models,
% legislative delegation, administrative discretion, reproducible
% workflows, quantitative legal coding


% ------------------------------------------------------------
% INTRODUCTION, PAGE 1
% ------------------------------------------------------------
% 1. The current statement that one-shot prompt classification is
%    "insufficient" is too absolute.
%
% REPLACE WITH:

Although one-shot LLM classification provides a simple baseline, it
can obscure the intermediate judgments underlying a final rank and can
produce outputs that are inconsistent with the measurement model.

% 2. Replace "Standard prompt-engineering approaches also remain
%    difficult to audit" with the more precise phrase below.

Single-prompt approaches are also difficult to audit because the
relationships among delegation, authority, constraints, and the final
classification remain embedded in one generative output.

% 3. Replace language saying that the prototype "renders" prompts and
%    settings as a graph.
%
% USE:

We implement the approach in Delegation Workflow Studio, a research
prototype that represents and executes prompts, settings, intermediate
outputs, and validation rules as a versioned and auditable workflow
graph.

% 4. The dashboard footnote appears twice on page 1. Retain it once
%    only. For the double-blind submission, remove or anonymize the
%    dashboard link, author names, affiliations, email addresses, and
%    identifying PDF metadata.
%
% 5. Replace the existing contribution list with the contribution list
%    provided earlier in this handoff. Replicability must appear first.


% ------------------------------------------------------------
% SECTION II.C, THEORY-GUIDED COMPUTATIONAL INFERENCE, PAGE 2
% ------------------------------------------------------------
% The current arrow chain implies that authority is evaluated before
% constraints and that the two are not parallel components.
%
% REPLACE:
% Delegation -> Authority -> Constraints -> Residual discretion
%
% WITH:

% Delegation -> {Authority, Constraints} -> Discretion rank

% Suggested prose:

Delegation is logically prior to discretion. Conditional on a
qualifying delegation, the workflow evaluates authority and constraints
as separate components and then uses their relationship to interpret or
validate the final discretion rank.


% ------------------------------------------------------------
% SECTION II.D, PROBLEM FORMULATION, PAGE 2
% ------------------------------------------------------------
% Replace this subsection in full with the formal-model section earlier
% in this handoff. This resolves the conflict between the deterministic
% theoretical mapping and the three rank-classification strategies.


% ------------------------------------------------------------
% SECTION IV AND TABLE I, PAGES 2--3
% ------------------------------------------------------------
% 1. Insert the revised architecture diagram supplied in this handoff.
%    It must visually separate:
%      (a) the deterministic theory-based mapping;
%      (b) the selected rank-classification strategy; and
%      (c) the consistency comparison.
%
% 2. Replace Table I with the revised table supplied earlier.
%
% 3. Replace "the workflow makes the measurement assumptions
%    falsifiable" with:

The workflow makes the measurement assumptions inspectable and
testable because researchers can locate a disagreement at the
Delegation Gate, authority classification, constraint classification,
rank-classification strategy, or consistency check.

% "Falsifiable" is too strong and is not the right description of an
% execution trace.
%
% 4. In the Validation-node definition, replace "throw errors" with:
%    "record validation flags" or "raise validation errors," depending
%    on what the software actually does.
%
% 5. Use "rank-classification strategy" everywhere. Do not use
%    "ranking head," which can imply a trained neural-network layer.


% ------------------------------------------------------------
% SECTION VI.A--B, BENCHMARK AND EVALUATION LEVELS, PAGES 4--5
% ------------------------------------------------------------
% 1. Audit and reconcile every sample count before submission.
%    Confirm separately:
%      - number of laws with delegation labels;
%      - number of delegation-positive laws;
%      - number of laws with complete expert-coded discretion ranks;
%      - number actually run for each model and strategy;
%      - number retained after failures.
%
% 2. Do not call the 15-law sample a test set. It is a development set
%    because it was used iteratively to revise the workflow.
%
% 3. Replace "full available discretionary corpus" throughout with:
%    "expert-coded discretion sample" or "delegated-law subset with
%    complete discretion labels," whichever is accurate.
%
% 4. Use N=169 only for analyses that actually use all 169 laws. Do not
%    label rank-classification tables N=169 if Rank 0 cases are excluded
%    or discretion labels are missing.


% ------------------------------------------------------------
% SECTION VI.D, EXPERIMENTAL PROTOCOL, PAGE 5
% ------------------------------------------------------------
% 1. Verify every provider name and exact API model identifier against
%    the execution logs. Report stable API identifiers, not only product
%    or marketing names.
%
% 2. Replace the temperature statement. Temperature zero reduces
%    sampling variation but does not guarantee identical outputs.
%
% REPLACE WITH:

The system uses a temperature of 0.0 to reduce sampling variation.
Because commercial APIs may still exhibit nondeterminism or provider-
side changes, output stability is evaluated through repeated runs and
the exact model identifiers and execution timestamps are retained.

% 3. Replace "If a node fails to compile or times out" with:

If a model call fails, returns an invalid structured output, or times
out, the system retries the call up to three times. Persistent failures
are recorded in the execution log and excluded or reported according to
the stated evaluation protocol.

% 4. State whether retries reuse the identical prompt and settings.
%    State how failed laws affect N in each table.


% ------------------------------------------------------------
% SECTION VI.E, METRICS, PAGE 5
% ------------------------------------------------------------
% The current binary-metric sentence repeats recall measures.
%
% REPLACE WITH:

For the binary delegation task, we report accuracy, balanced accuracy,
class-specific precision and recall, macro-F1, and the confusion matrix.
For the ordinal discretion task, we report exact-match accuracy,
within-one accuracy, mean absolute error, linear weighted kappa,
quadratic weighted kappa, and the ordinal confusion matrix.

% Use distinct notation for linear and quadratic weighted kappa in the
% equations or table headings if notation is shown.


% ------------------------------------------------------------
% SECTION VII AND TABLES IV--VIII, PAGES 5--7
% ------------------------------------------------------------
% 1. Define the comparison system currently called "Looser Screening."
%    The label is informal and unclear. Use a descriptive name based on
%    the actual prompt design, such as "Broad Delegation Screen" or
%    "High-Recall Delegation Screen," and define it once. Do not rename
%    it unless that description accurately reflects the implementation.
%
% 2. Table V: define "Logical Consistency" explicitly as agreement
%    between the strategy-specific rank and the rank implied by the
%    component classifications. Define "Invalid Output Rate" as the
%    percentage of runs that fail schema validation after the permitted
%    retries.
%
% 3. Table VI: use the actual rank-evaluation N for every model and
%    strategy. If N differs because of failures, show N in the table or
%    footnote rather than presenting one universal N.
%
% 4. Table VII: define repeat-run agreement. Recommended definition:
%    the share of cases for which all three runs return the same value.
%    State whether agreement is calculated only among completed runs.
%
% Suggested definition:

% \[
% \mathrm{Agreement}(z)=\frac{1}{N}\sum_{i=1}^{N}
% \mathbf{1}\left[z_i^{(1)}=z_i^{(2)}=z_i^{(3)}\right].
% \]

% 5. Table VIII: define runtime, token, cost, completion, and retry
%    measures in the text. State the currency and date of API pricing.
%
% 6. Remove every "TBD" before submission. If a planned analysis cannot
%    be completed, remove the table and narrow the empirical claim rather
%    than leaving an empty result structure.


% ------------------------------------------------------------
% SECTION VII.B, DIAGNOSING DISAGREEMENT, PAGE 5
% ------------------------------------------------------------
% The PL83-577 example is useful. Clarify that it is a diagnostic case,
% not evidence that the historical benchmark is wrong.
%
% SUGGESTED FINAL SENTENCE:

The trace does not by itself determine which label is correct; it makes
the source of the disagreement visible for expert adjudication.


% ------------------------------------------------------------
% SECTION VII.G, LIMITATIONS, PAGE 5
% ------------------------------------------------------------
% Replace "raw statutory text blocks" with "full statutory text."
%
% Replace "the modular workflow enforcing hard logical constraints"
% with:

Although deterministic validation enforces specified consistency rules,
the underlying LLM classifications remain sensitive to model choice,
prompt wording, nondeterministic API behavior, and provider-side model
updates.


% ------------------------------------------------------------
% SECTION VII.H, NOVELTY CLAIM, PAGE 5
% ------------------------------------------------------------
% Retain this subsection, but use the shorter novelty statement supplied
% earlier in this handoff. Do not frame human-in-the-loop annotation or
% theory-guided AI as independent new contributions.


% ------------------------------------------------------------
% SECTION VIII, REPLICATION WORKFLOW, PAGE 6
% ------------------------------------------------------------
% 1. Keep this section, but shorten it if space is needed.
%
% 2. Replace product-specific wording such as:
%    "Linking a campaign copies the workflow's definition and revision
%    directly to the dashboard."
%
% USE:

Before an evaluation run, the workflow definition, prompts, model
settings, and source set are frozen and assigned a version identifier.
Subsequent edits create a new workflow version and do not alter the
record of prior runs.

% 3. In Step 7, replace "rerun the fixed development set" with wording
%    that distinguishes development and evaluation data:

If a disagreement reveals a general coding-rule problem, revise the
workflow as a new version, rerun the development set, and evaluate the
revised version on data not used to make that change when such data are
available.


% ------------------------------------------------------------
% SECTION IX, FUTURE WORK, PAGE 6
% ------------------------------------------------------------
% Remove internal prompt-version labels (v8, v8.2, v8.4, and v9) from
% the main paper. They are appropriate for the replication materials,
% not for the conference narrative.
%
% Also remove any future-work sentence promising analyses that the
% abstract or results section already claims have been completed.
%
% REPLACEMENT PARAGRAPH:

Future work will evaluate the workflow on a held-out set of laws,
measure component-level authority and constraint performance, and test
external validity in an additional policy domain. Further system work
will strengthen workflow-version comparison, approval controls, and
case-adjudication interfaces. Independent expert recoding would also
permit direct estimation of intercoder reliability and clarify the
upper bound on model-benchmark agreement.


% ------------------------------------------------------------
% SECTION X, CONCLUSION, PAGES 6--7
% ------------------------------------------------------------
% Retain the final sentence that the system makes judgment explicit,
% structured, inspectable, and reproducible. Precede it with the
% positioning language supplied earlier:

The resulting system is a replicable measurement instrument whose
computational structure is guided by an established theory of
delegation and discretion. Replicability is the primary contribution;
theory-guided design explains why the workflow's stages, dependencies,
and validation rules are substantively meaningful.


% ------------------------------------------------------------
% REFERENCES, PAGES 7--8
% ------------------------------------------------------------
% 1. Remove every phrase saying "bibliographic details to be verified."
%    Verify and complete each citation before submission.
%
% 2. Check author lists, titles, venue names, years, page ranges, and DOI
%    information against the original publication records.
%
% 3. Ensure that self-citations do not reveal authorship in a double-
%    blind submission through phrases such as "our prior work." Cite
%    them in the third person in the manuscript.


% ============================================================
% PRIORITY B: GLOBAL TERMINOLOGY CONSISTENCY
% ============================================================
%
% Use consistently:
%   Delegation Gate
%   authority classification
%   constraint classification
%   rank-classification strategy
%   Sequential Ordinal Cascade
%   Direct Multiclass Rank Prediction
%   Two-Band Rank Classification
%   theory-based mapping or theoretical mapping
%   rank implied by the component classifications
%   consistency indicator or inconsistency flag
%   execution trace
%   workflow version
%   expert-coded benchmark
%
% Avoid or replace:
%   ranking head
%   strict reproducibility
%   ensures identical output
%   fails to compile
%   falsifiable workflow
%   looser screening (unless explicitly defined)
%   full available discretionary corpus
%   raw statutory text blocks
%   admissible rank relationships
%   internal-consistency failure
%   strategy acronyms unless the paper genuinely needs them
%
% Use "replicable" for the measurement instrument and "reproducible"
% for the preserved computational execution record. Define this usage
% once if both terms remain in the manuscript.
% ============================================================
